\documentclass[../../thesis_main.tex]{subfiles}
\graphicspath{{\subfix{../../figures/}}}

\begin{document}

\chapter*{Conclusions}
\addcontentsline{toc}{chapter}{Conclusions}
\markboth{Conclusions}{}

This thesis investigated the experimental and numerical components required for the preparation, optical trapping, transport, and cooling of cold $^{87}$Rb atoms in the context of a Hollow-Core Photonic Crystal Fiber experiment.

The preparation of the atomic sample was optimized by adjusting the optical-molasses parameters and compensation magnetic fields, yielding a minimum temperature of $T = (8.4 \pm 0.2)\,\mu\text{K}$ measured via ballistic expansion. The loading of the Optical Dipole Trap was then optimized by translating the MOT to improve its spatial overlap with the ODT beam. This resulted in a measured loading efficiency of $\eta = (0.161 \pm 0.005)\,\%$ and a \CCn{conservatively estimated trapped population}{``Stima conservativa'' e barra d'errore $\pm0.1$ dicono due cose diverse: la prima e' un limite inferiore, la seconda un'incertezza statistica. Conviene dire cosa rappresenta il $\pm$ (dispersione della calibrazione in fluorescenza) e separatamente che il valore centrale e' un limite inferiore, per i due motivi che il capitolo 3 elenca: i 600~$\mu$s di ODT spento e la calibrazione a due livelli.} of $N_{\mathrm{ODT}} = (1.0 \pm 0.1)\times 10^{5}\ \text{atoms}$.

To establish the optical geometry for a future conveyor-belt implementation, an injection line for a 1064~nm polarization-maintaining single-mode fiber was designed, characterized, and spatially aligned with the counter-propagating field delivered through the HCPCF. The stability of the field transmitted through this fiber was characterized, with the fractional power and angular Allan deviations reaching minima of $8\times10^{-4}$ at $0.1$~s and $7\times10^{-5}$~rad at $3$~s, respectively. At longer averaging times, both measurements were increasingly affected by low-frequency technical noise and environmental drifts. While these measurements demonstrate the stability of the injected branch, characterizing the complete conveyor-belt potential will require an independent analysis of the HCPCF branch.

The effect of optical pumping on EIT cooling was investigated numerically, moving from an ideal three-level $\Lambda$ model to a 24-level system including the hyperfine and Zeeman structure of the $^{87}$Rb $D_{2}$ transition. \CCerr{The multilevel simulations reveal that spontaneous emission into uncoupled ground states proressively interrupts the cooling cycle.}{Due cose introdotte dalla revisione di \texttt{a3a1f66}. Refuso: ``proressively'' $\to$ ``progressively''. E soprattutto e' stata CANCELLATA la frase ``These results identify repumping and internal-state management as essential requirements, while noting that cooling dynamics are substantially slower in the multilevel system than in the ideal model''. Era l'unico punto di tutta la tesi, fuori dal capitolo 3, in cui compariva la riserva sulla lentezza della dinamica multilivello: l'abstract non l'ha mai avuta, e ora non ce l'hanno piu' nemmeno le conclusioni. Con il modello che per giunta trascura il rinculo dell'emissione spontanea (vedi la nota nel capitolo 2), toglierla lascia l'affermazione sul ground state senza nessuna qualificazione. La rimetterei.} Including an active repump field sustains population recycling, \CCn{allowing the simulated mean phonon number to approach the motional ground state}{Anche qui manca il valore di $\langle n\rangle$. Aggiungerei il numero e il confronto con il limite $(\gamma/4\Delta)^2\simeq3.7\times10^{-4}$ atteso per $\Delta_c=13\gamma$: e' l'unico modo per far capire se il ``ground state'' e' raggiunto davvero o solo avvicinato.}. Future simulations will still need to incorporate specific perturbations expected in the fiber environment, such as position-dependent light shifts and local magnetic-field gradients.

Overall, this work establishes the main experimental infrastructure and numerical framework needed for controlled atom transport into the HCPCF. The next steps are the experimental realization of the conveyor-belt motion, the characterization of transport, and the implementation of EIT cooling under the experimental optical and magnetic conditions inside the hollow core.

\end{document}